Durant una tempesta cau un llamp pel qual circula un corrent elèctric de 400~kA. Suposeu que la intensitat del corrent del llamp és constant durant els 50~$\mu$s que dura.

\begin{apartat}{1,25}
Quina és la càrrega elèctrica total que ha transportat aquest llamp? Quin és el camp magnètic que crea aquest corrent a una distància de 100~m?
\end{apartat}

\begin{apartat}{1,25}
Quina força magnètica actua sobre una partícula carregada que es troba en repòs a aquesta mateixa distància? Justifiqueu la resposta.
\end{apartat}

\dades{%
  $\mu_0 = 4\pi\times 10^{-7}\ \mathrm{T\,m\,A^{-1}}$.\\
  $|e| = 1,602\times 10^{-19}\ \mathrm{C}$.\\[2pt]
  La intensitat del camp magnètic creat per un corrent rectilini és $B = \dfrac{\mu_0 I}{2\pi r}$, en què $r$ és la distància al corrent.}
